% TOP_PROBLEM: {"id": 149, "title": "N Queens Problem Asymptotics", "category_id": 2, "category": {"id": 2, "name": "combinatorics", "display_name": "Combinatorics", "description": "Counting problems, graph theory, discrete structures.", "slug": "combinatorics", "order_index": 2, "created_at": "2026-07-31T15:26:25.671Z"}, "status": "open"}
% FIRST_POSTED: null
% ENTRY_KIND: statement_only
% Imported from: batch03.tex; UnsolvedMath ID 149
\documentclass[11pt]{article}
\usepackage[margin=25mm]{geometry}
\usepackage{fontspec}
\setmainfont{Latin Modern Roman}
\usepackage{amsmath,amssymb,mathtools}
\usepackage{unicode-math}
\setmathfont{Latin Modern Math}
\usepackage{xurl,hyperref}
\hypersetup{colorlinks=true,urlcolor=blue}
\setcounter{secnumdepth}{0}
\setlength{\parindent}{0pt}
\setlength{\parskip}{5pt}
\setlength{\emergencystretch}{3em}
\newcommand{\sref}[2]{\hyperlink{source:#1:#2}{[S#2]}}
\newcommand{\cataloguescope}[1]{\paragraph{Recorded target.}#1}
\begin{document}
% UnsolvedMath ID 149; source catalogue code GREEN-061
\setcounter{section}{148}
\section{N Queens Problem Asymptotics}\label{149}
{\small\sffamily UnsolvedMath ID 149 · Combinatorics}
\subsection{Definitions and mathematical statement}

\paragraph{Shared notation.}$\mathbb N=\{1,2,\ldots\}$, and $\mathbb Z,\mathbb Q,\mathbb R,\mathbb C$ denote the integers, rationals, reals and complex numbers. Any other conventions are specified locally.


Write $[n]=\{1,\ldots,n\}$, and let $S_n$ be the permutations of $[n]$. A configuration of $n$ queens on the ordinary $n\times n$ chessboard has exactly one queen in each row and column, hence is represented uniquely by a permutation $\pi\in S_n$, with a queen at $(i,\pi(i))$. Nonattack along the two diagonal directions is the requirement
\[
 |\pi(i)-\pi(j)|\ne |i-j|\qquad(1\leq i<j\leq n).
\]
Define
\[
 Q(n)=\bigl|\{\pi\in S_n:|\pi(i)-\pi(j)|\ne|i-j|\text{ for every }i<j\}\bigr|.
\]
The board has no toroidal identifications, and configurations related by a rotation or reflection are counted separately.
\paragraph{Statement recorded in the supplied source.}
Determine the asymptotic growth of $Q(n)$ as $n\to\infty$. In particular, determine the limit of $Q(n)^{1/n}/n$, or equivalently the constant $\alpha$ in the exponential-scale form
\[
 Q(n)=n^n\exp\bigl(-\alpha n+o(n)\bigr).
\] \sref{149}{2}
\subsection{Short English statement}
Determine the asymptotic number of placements of n mutually nonattacking queens on an ordinary n-by-n board.
\subsection{Sources}
\begin{description}
\item[\hypertarget{source:149:1}{S1}] ULAM.AI and the UnsolvedMath Contributors, \emph{UnsolvedMath: A Curated Collection of Open Mathematics Problems}, record 149 (GREEN-061). CC BY 4.0. \url{https://huggingface.co/datasets/ulamai/UnsolvedMath}.
\item[\hypertarget{source:149:2}{S2}] Michael Simkin, \emph{The number of $n$-queens configurations}, arXiv:2107.13460v3 (2022), §1, Theorem 1.1. This theorem answers the recorded exponential-growth question: $Q(n)^{1/n}/n\to e^{-\alpha}$ for a variationally characterized constant $1.94<\alpha<1.9449$. \url{https://arxiv.org/abs/2107.13460}.
\item[\hypertarget{source:149:3}{S3}] Ben Green, \emph{100 Open Problems}, maintained author list, Problem 97 and its comments, accessed 7 October 2026. \url{https://people.maths.ox.ac.uk/greenbj/papers/open-problems.pdf}.
\end{description}
\label{end:149}
\end{document}
